\begin{table}
\begin{tcolorbox}[tabularx={X|p{2cm}|X|},enhanced,fonttitle=\bfseries\large,fontupper=\normalsize\sffamily,
colback=yellow!10!white,colframe=red!50!black,colbacktitle=blue!30!white,
coltitle=black,title=Options for CRC-Codes]
\hline
Command & Arguments & Purpose \\
\hline
\hline
\verb'-input'  &  fname  &  Input file name. \\
\hline
\verb'-output'  &  fname  &  Output file name. \\
\hline
\verb'-crc_type'  &  type  &  The type of CRC-code. \\
\hline
\verb'-block_length'  &  $L$  &  Set the block length to $L$ field elements. \\
\hline
\verb'-block_based_error_generator'  &    &  Apply block-based error generator. \\
\hline
\verb'-file_based_error_generator'  &  threshold  &  Apply file-based error generator. \\
\hline
\verb'-nb_repeats'  &  $N$  &  Set the number of repeats to $N$. \\
\hline
\verb'-threshold'  &  $t$  &  Set probability of error per experiment to $t/1000000$. \\
\hline
\verb'-error_log'  &  fname  &  Set file name for error logging. \\
\hline
\verb'-selected_block'  &  $i$  &  Set block number. \\
\hline
\end{tcolorbox}
\caption{\label{tab:CRC:options}Options for CRC-Codes}
\end{table}
